% ============================================================
\chapter{Probability Review}
\label{ch:review}
% ============================================================

Before diving into the study of stochastic processes --- those random
functions that evolve over time --- we must ensure that our probabilistic
foundations are solid. This chapter reviews the essential tools: probability
spaces, random variables, conditional expectation, and modes of convergence.
Readers already comfortable with these notions may skim this chapter as a
reference; otherwise, this is where everything begins.

% ------------------------------------------------------------
\section{Probability Spaces}
% ------------------------------------------------------------

Let us recall the three fundamental ingredients of Kolmogorov's framework.

\begin{definition}[$\sigma$-Algebra]
Let $\Omega$ be a non-empty set. A \emph{$\sigma$-algebra} (or $\sigma$-field)
$\mathcal{F}$ on $\Omega$ is a collection of subsets of $\Omega$ satisfying:
\begin{enumerate}
  \item $\Omega \in \mathcal{F}$,
  \item if $A \in \mathcal{F}$, then $A^c \in \mathcal{F}$ (closed under complementation),
  \item if $(A_n)_{n \geq 1} \subset \mathcal{F}$, then $\bigcup_{n=1}^{\infty} A_n \in \mathcal{F}$
    (closed under countable unions).
\end{enumerate}
The pair $(\Omega, \mathcal{F})$ is called a \emph{measurable space}.
\end{definition}

\begin{definition}[Probability Space]
A \emph{probability space} is a triple $(\Omega, \mathcal{F}, \PP)$ where:
\begin{itemize}
  \item $\Omega$ is the \emph{sample space} (set of all possible outcomes),
  \item $\mathcal{F}$ is a $\sigma$-algebra on $\Omega$ (the events),
  \item $\PP : \mathcal{F} \to [0,1]$ is a \emph{probability measure}, i.e.\ a
    $\sigma$-additive measure with $\PP(\Omega) = 1$.
\end{itemize}
\end{definition}

\begin{remark}
The $\sigma$-additivity of $\PP$ means: for any sequence $(A_n)_{n \geq 1}$
of pairwise disjoint events,
\[
  \PP\Bigl(\bigcup_{n=1}^{\infty} A_n\Bigr) = \sum_{n=1}^{\infty} \PP(A_n).
\]
\end{remark}

\begin{definition}[Generated $\sigma$-Algebra]
Let $\mathcal{C}$ be a collection of subsets of $\Omega$. The \emph{$\sigma$-algebra
generated} by $\mathcal{C}$, denoted $\sigma(\mathcal{C})$, is the smallest
$\sigma$-algebra containing $\mathcal{C}$:
\[
  \sigma(\mathcal{C}) = \bigcap \bigl\{ \mathcal{F} : \mathcal{F} \text{ is a }
  \sigma\text{-algebra on } \Omega,\; \mathcal{C} \subset \mathcal{F} \bigr\}.
\]
In particular, the \emph{Borel $\sigma$-algebra} $\mathcal{B}(\R)$ is generated
by the open sets of $\R$.
\end{definition}

\section{Random Variables}

\begin{definition}[Random Variable]
Let $(\Omega, \mathcal{F}, \PP)$ be a probability space. A \emph{random variable}
(r.v.) is a measurable map $X : (\Omega, \mathcal{F}) \to (\R, \mathcal{B}(\R))$,
i.e.:
\[
  \forall B \in \mathcal{B}(\R), \quad X^{-1}(B) = \{\omega \in \Omega : X(\omega) \in B\} \in \mathcal{F}.
\]
More generally, a \emph{random vector} is a measurable map
$X : (\Omega, \mathcal{F}) \to (\R^d, \mathcal{B}(\R^d))$.
\end{definition}

\begin{definition}[$\sigma$-Algebra Generated by a Random Variable]
The $\sigma$-algebra generated by a random variable $X$ is
\[
  \sigma(X) = \{X^{-1}(B) : B \in \mathcal{B}(\R)\}.
\]
This is the smallest $\sigma$-algebra making $X$ measurable. Intuitively,
$\sigma(X)$ represents the \emph{information} contained in the observation of $X$.
\end{definition}

\begin{theorem}[Doob's Measurability Lemma]
\label{thm:doob_lemma}
Let $X$ be a random variable and $\mathcal{G} = \sigma(X)$. A random variable
$Y$ is $\mathcal{G}$-measurable if and only if there exists a Borel measurable
function $g : \R \to \R$ such that $Y = g(X)$.
\end{theorem}

\section{Expectation and Integration}

\begin{definition}[Expectation]
The \emph{expectation} of a non-negative random variable $X \geq 0$ is
\[
  \E[X] = \int_{\Omega} X(\omega) \, d\PP(\omega).
\]
For a general r.v., $\E[X] = \E[X^+] - \E[X^-]$ when at least one of these
quantities is finite, where $X^+ = \max(X, 0)$ and $X^- = \max(-X, 0)$.
We say $X$ is \emph{integrable} if $\E[\abs{X}] < \infty$.
\end{definition}

\begin{theorem}[Lebesgue Dominated Convergence]
\label{thm:dom_conv}
Let $(X_n)_{n \geq 1}$ be a sequence of r.v.\ converging a.s.\ to $X$. If there
exists an integrable r.v.\ $Y$ such that $\abs{X_n} \leq Y$ a.s.\ for all $n$,
then $X$ is integrable and
\[
  \lim_{n \to \infty} \E[X_n] = \E[X].
\]
\end{theorem}

\begin{theorem}[Monotone Convergence]
\label{thm:mon_conv}
Let $(X_n)_{n \geq 1}$ be an increasing sequence of non-negative r.v.:
$0 \leq X_1 \leq X_2 \leq \cdots$ a.s. Then
\[
  \lim_{n \to \infty} \E[X_n] = \E\bigl[\lim_{n \to \infty} X_n\bigr].
\]
\end{theorem}

\begin{theorem}[Fatou's Lemma]
Let $(X_n)_{n \geq 1}$ be a sequence of non-negative r.v. Then
\[
  \E\bigl[\liminf_{n \to \infty} X_n\bigr] \leq \liminf_{n \to \infty} \E[X_n].
\]
\end{theorem}

\section{$L^p$ Spaces and Inequalities}

\begin{definition}[$L^p$ Space]
For $p \in [1, \infty)$, the space $L^p(\Omega, \mathcal{F}, \PP)$ is the set
of (equivalence classes of) r.v.\ $X$ with $\E[\abs{X}^p] < \infty$, equipped
with the norm
\[
  \norm{X}_{L^p} = \bigl(\E[\abs{X}^p]\bigr)^{1/p}.
\]
The space $L^{\infty}$ is the set of essentially bounded r.v.
\end{definition}

\begin{theorem}[H\"older's Inequality]
Let $p, q \in (1, \infty)$ with $\frac{1}{p} + \frac{1}{q} = 1$. For
$X \in L^p$ and $Y \in L^q$,
\[
  \E[\abs{XY}] \leq \norm{X}_{L^p} \norm{Y}_{L^q}.
\]
The case $p = q = 2$ gives the \emph{Cauchy--Schwarz inequality}.
\end{theorem}

\begin{theorem}[Jensen's Inequality]
\label{thm:jensen}
If $X$ is integrable and $\varphi : \R \to \R$ is convex, then
\[
  \varphi(\E[X]) \leq \E[\varphi(X)],
\]
provided $\E[\abs{\varphi(X)}] < \infty$.
\end{theorem}

\section{Conditional Expectation}

\begin{definition}[Conditional Expectation]
\label{def:cond_exp}
Let $X \in L^1(\Omega, \mathcal{F}, \PP)$ and let $\mathcal{G} \subset \mathcal{F}$
be a sub-$\sigma$-algebra. The \emph{conditional expectation} of $X$ given
$\mathcal{G}$, denoted $\E[X \mid \mathcal{G}]$, is the unique (a.s.) random
variable satisfying:
\begin{enumerate}
  \item $\E[X \mid \mathcal{G}]$ is $\mathcal{G}$-measurable,
  \item for all $G \in \mathcal{G}$, $\displaystyle \int_G \E[X \mid \mathcal{G}] \, d\PP
    = \int_G X \, d\PP$.
\end{enumerate}
Existence is guaranteed by the Radon--Nikodym theorem.
\end{definition}

\begin{theorem}[Properties of Conditional Expectation]
\label{thm:ce_props}
Let $X, Y \in L^1$ and let $\mathcal{G}, \mathcal{H}$ be sub-$\sigma$-algebras of $\mathcal{F}$.
\begin{enumerate}
  \item \textbf{Linearity}: $\E[aX + bY \mid \mathcal{G}] = a\E[X \mid \mathcal{G}]
    + b\E[Y \mid \mathcal{G}]$ a.s.
  \item \textbf{Positivity}: if $X \geq 0$ a.s., then $\E[X \mid \mathcal{G}] \geq 0$ a.s.
  \item \textbf{Tower property (smoothing)}: if $\mathcal{H} \subset \mathcal{G}$,
    then $\E[\E[X \mid \mathcal{G}] \mid \mathcal{H}] = \E[X \mid \mathcal{H}]$ a.s.
  \item \textbf{Pull-out property}: if $Y$ is $\mathcal{G}$-measurable and bounded,
    $\E[XY \mid \mathcal{G}] = Y \E[X \mid \mathcal{G}]$ a.s.
  \item \textbf{Independence}: if $\sigma(X)$ and $\mathcal{G}$ are independent,
    $\E[X \mid \mathcal{G}] = \E[X]$ a.s.
  \item \textbf{Conditional Jensen}: if $\varphi$ is convex and $\varphi(X) \in L^1$,
    $\varphi(\E[X \mid \mathcal{G}]) \leq \E[\varphi(X) \mid \mathcal{G}]$ a.s.
\end{enumerate}
\end{theorem}

\begin{intuition}[title=\textbf{Geometric Interpretation}]
In $L^2(\Omega, \mathcal{F}, \PP)$, the conditional expectation
$\E[X \mid \mathcal{G}]$ is the \emph{orthogonal projection} of $X$ onto the
closed subspace $L^2(\Omega, \mathcal{G}, \PP)$. This justifies:
\[
  \E\bigl[\bigl(X - \E[X \mid \mathcal{G}]\bigr)^2\bigr]
  \leq \E\bigl[(X - Z)^2\bigr]
\]
for all $Z \in L^2(\Omega, \mathcal{G}, \PP)$.
\end{intuition}

\section{Independence}

\begin{definition}[Independence of $\sigma$-Algebras]
Sub-$\sigma$-algebras $\mathcal{G}_1, \ldots, \mathcal{G}_n$ of $\mathcal{F}$
are \emph{independent} if for every choice $A_i \in \mathcal{G}_i$:
\[
  \PP(A_1 \cap \cdots \cap A_n) = \PP(A_1) \cdots \PP(A_n).
\]
Random variables $X_1, \ldots, X_n$ are independent if $\sigma(X_1), \ldots,
\sigma(X_n)$ are independent.
\end{definition}

\begin{definition}[Mutual Independence]
A family $(X_i)_{i \in I}$ of random variables is \emph{mutually independent}
if for every finite subset $J \subset I$, the random variables $(X_j)_{j \in J}$
are independent.
\end{definition}

\begin{theorem}[Kolmogorov's Zero--One Law]
\label{thm:kolmo01}
Let $(X_n)_{n \geq 1}$ be a sequence of independent r.v. The \emph{tail
$\sigma$-algebra}
\[
  \mathcal{T} = \bigcap_{n=1}^{\infty} \sigma(X_n, X_{n+1}, \ldots)
\]
is \emph{trivial}: for every $A \in \mathcal{T}$, $\PP(A) \in \{0, 1\}$.
\end{theorem}

\section{Modes of Convergence}

\begin{definition}[Types of Convergence]
Let $(X_n)_{n \geq 1}$ be a sequence of r.v.\ and $X$ a r.v.\ defined on
$(\Omega, \mathcal{F}, \PP)$.
\begin{enumerate}
  \item \textbf{Almost sure convergence}: $X_n \xrightarrow{\text{a.s.}} X$
    if $\PP(\{\omega : X_n(\omega) \to X(\omega)\}) = 1$.
  \item \textbf{$L^p$ convergence}: $X_n \xrightarrow{L^p} X$ if
    $\E[\abs{X_n - X}^p] \to 0$.
  \item \textbf{Convergence in probability}: $X_n \xrightarrow{\PP} X$ if
    for every $\varepsilon > 0$, $\PP(\abs{X_n - X} > \varepsilon) \to 0$.
  \item \textbf{Convergence in distribution}: $X_n \xrightarrow{\mathcal{L}} X$ if
    $\E[f(X_n)] \to \E[f(X)]$ for every bounded continuous $f : \R \to \R$.
\end{enumerate}
\end{definition}

\begin{keyformulas}[title=\textbf{Hierarchy of Convergences}]
\[
  \text{a.s.} \implies \text{in probability} \implies \text{in distribution}
\]
\[
  L^p \implies L^q \;\;(q \leq p) \implies \text{in probability} \implies \text{in distribution}
\]
The converses are false in general. However:
\begin{itemize}
  \item $X_n \xrightarrow{\PP} X$ $\implies$ there exists a subsequence $X_{n_k} \xrightarrow{\text{a.s.}} X$.
  \item If $X_n \xrightarrow{\mathcal{L}} c$ (constant), then $X_n \xrightarrow{\PP} c$.
\end{itemize}
\end{keyformulas}

\section{Laws of Large Numbers}

\begin{theorem}[Strong Law of Large Numbers]
\label{thm:slln}
Let $(X_n)_{n \geq 1}$ be a sequence of i.i.d.\ integrable r.v.\ with mean
$\mu = \E[X_1]$. Then
\[
  \frac{1}{n} \sum_{k=1}^{n} X_k \xrightarrow{\text{a.s.}} \mu.
\]
\end{theorem}

\begin{theorem}[Central Limit Theorem]
\label{thm:clt}
Let $(X_n)_{n \geq 1}$ be i.i.d.\ with $\E[X_1] = \mu$ and
$\Var(X_1) = \sigma^2 \in (0, \infty)$. Then
\[
  \frac{1}{\sqrt{n}} \sum_{k=1}^{n} \frac{X_k - \mu}{\sigma}
  \xrightarrow{\mathcal{L}} \mathcal{N}(0, 1).
\]
\end{theorem}

\section{Borel--Cantelli Lemma}

\begin{theorem}[Borel--Cantelli Lemma]
\label{thm:borel_cantelli}
Let $(A_n)_{n \geq 1}$ be a sequence of events.
\begin{enumerate}
  \item If $\sum_{n=1}^{\infty} \PP(A_n) < \infty$, then
    $\PP(\limsup_{n} A_n) = 0$, i.e.\ a.s.\ only finitely many $A_n$ occur.
  \item If the $A_n$ are \emph{independent} and $\sum_{n=1}^{\infty} \PP(A_n) = \infty$,
    then $\PP(\limsup_{n} A_n) = 1$, i.e.\ a.s.\ infinitely many $A_n$ occur.
\end{enumerate}
\end{theorem}

\begin{warning}[title=\textbf{Independence Hypothesis}]
The second part of the Borel--Cantelli lemma requires independence of the events.
Without this hypothesis, divergence of $\sum \PP(A_n)$ does not guarantee
$\PP(\limsup A_n) = 1$.
\end{warning}

\section{Characteristic Functions}

\begin{definition}[Characteristic Function]
The \emph{characteristic function} of a r.v.\ $X$ is
\[
  \varphi_X(t) = \E[e^{itX}], \quad t \in \R.
\]
\end{definition}

\begin{proposition}[Properties]
\begin{enumerate}
  \item $\varphi_X(0) = 1$ and $\abs{\varphi_X(t)} \leq 1$ for all $t$.
  \item $\varphi_X$ is uniformly continuous on $\R$.
  \item $\varphi_X$ characterises the law of $X$: if $\varphi_X = \varphi_Y$,
    then $X \overset{\mathcal{L}}{=} Y$.
  \item If $X$ and $Y$ are independent, $\varphi_{X+Y}(t) = \varphi_X(t) \varphi_Y(t)$.
  \item If $\E[\abs{X}^n] < \infty$, then $\varphi_X^{(n)}(0) = i^n \E[X^n]$.
\end{enumerate}
\end{proposition}

\begin{theorem}[L\'evy's Continuity Theorem]
\label{thm:levy}
$X_n \xrightarrow{\mathcal{L}} X$ if and only if $\varphi_{X_n}(t) \to \varphi_X(t)$
for all $t \in \R$.
\end{theorem}

\section{Exercises}

\begin{exercise}
Let $(\Omega, \mathcal{F}, \PP)$ be a probability space and $X$ a non-negative r.v.
Show that $\E[X] = \int_0^{\infty} \PP(X > t) \, dt$.
\end{exercise}

\begin{exercise}
Show that a.s.\ convergence does not imply $L^1$ convergence by constructing
a counterexample.
\end{exercise}

\begin{exercise}
Let $(X_n)_{n \geq 1}$ be independent r.v.\ with
$\PP(X_n = n) = 1/n$ and $\PP(X_n = 0) = 1 - 1/n$. Show that $X_n \to 0$ a.s.
\end{exercise}

\begin{exercise}
Prove the tower property of conditional expectation
(Theorem~\ref{thm:ce_props}, item~3).
\end{exercise}

\begin{exercise}
Let $X$ be integrable and $\mathcal{G}_n \uparrow \mathcal{G}_{\infty}$ an
increasing filtration of sub-$\sigma$-algebras. Show that
$\E[X \mid \mathcal{G}_n] \to \E[X \mid \mathcal{G}_{\infty}]$ a.s.\ and in $L^1$.
\emph{(This is L\'evy's martingale convergence theorem.)}
\end{exercise}
